\chapter{Aktywne uczenie automatów}

Rozdział ten ma na celu przedstawienie ogólnego kontekstu aktywnego uczenia automatów, na którym opiera się dalsza
część pracy. Najpierw omówiony zostanie klasyczny model aktywnego uczenia automatów DFA oraz jego rozszerzenie
o dodatkową podpowiedź w postaci systemu SRS. Następnie przedstawione zostaną najważniejsze modyfikacje
konieczne przy przejściu od automatów DFA do automatów DVPA. Na końcu rozdziału krótko opisano rozwiązania heurystyczne
wykorzystane w implementacji przedstawionej w tej pracy.

\section{Uczenie się automatów DFA z poradą}

\textbf{Model aktywnego uczenia} zakłada istnienie nauczyciela, który potrafi odpowiadać na określony zestaw zapytań dotyczących nieznanego
języka docelowego $L$. Algorytm konstruujący automat $H$, nazywany hipotezą, który docelowo powinien rozpoznawać język $L$, nazywany jest uczniem. 

W klasycznym modelu Angluin~\cite{Angluin} nauczyciel odpowiada na dwa typy zapytań:
\begin{itemize}
  \item \textbf{zapytanie o należenie (ang. membership query) $MQ(w)$} -- zapytanie o~to, czy dane słowo $w$ należy do języka $L$.
  \item \textbf{zapytanie o równoważność (ang. equivalence query) $EQ(H)$} --  zapytanie o to, czy skonstruowana hipoteza $H$ rozpoznaje język $L$.
\end{itemize}
W odpowiedzi na zapytanie o równoważność nauczyciel może ją potwierdzić, co kończy działanie algorytmu, albo zwrócić kontrprzykład,
czyli słowo, które należy do dokładnie jednego z języków $L$ oraz $L(H)$.

\textbf{Algorytm L$^*$}, zaproponowany przez Angluin~\cite{Angluin}, stanowi realizację powyższego modelu dla języków regularnych.
W tym podejściu uczeń utrzymuje dwa zbiory: zbiór selektorów $S$ reprezentujący stany automatu i zbiór słów testowych $C$.
Relacja $\sim^C_L$ jest zdefiniowana w następujący sposób:
$$s_1 \sim^C_L s_2 \Leftrightarrow \forall_{c\in C} (s_1c \in L \Leftrightarrow s_2c \in L)$$
W trakcie trwania algorytmu chcemy utrzymywać dwa niezmienniki dla tych zbiorów:

\begin{enumerate}
  \item \textbf{Rozdzielność}: różne słowa z $S$ należą do różnych klas abstrakcji relacji $\sim^C_L$
  \item \textbf{Domkniętość}: dla każdego słowa $s \in S$ i $a \in \Sigma$ istnieje $s' \in S$, takie że $sa \sim^C_L s'$
\end{enumerate}

Jeśli $C = \Sigma^*$, to relacja $\sim^C_L$ jest dokładnie relacją $\sim_L$ z Twierdzenia Myhilla-Nerode'a.
Gwarantuje to, że rozszerzając zbiory $S$ i $C$, a przy tym utrzymując powyższe własności będziemy w stanie skonstruować minimalny automat DFA
rozpoznający język $L$.

Algorytm zaczyna działanie z $S = \{\epsilon\}$ i $C = \emptyset$, które spełniają powyższe własności. Następnie powtarza w pętli procedurę:

\begin{enumerate}
    \item Zaczynając z rozdzielnymi $S$ i $C$ domknij je rozszerzając zbiór $S$ o dodatkowe selektory w postaci $sa$,
    jeśli dla danych $s \in S$ i $a \in \Sigma$ nie istnieje $s'$ spełniający własność domkniętości.
    Jeśli własność rozdzielności nie jest zachowana rozszerz zbiór $C$ o odpowiedni sufiks.

    \item Skonstruuj hipotezę $H$ na podstawie $S$ i $C$ korzystając z własności domkniętości, następnie zadaj zapytanie $EQ(H)$.
    Jeśli nauczyciel potwierdzi równoważność zakończ działanie.

    \item Jeśli nauczyciel zwrócił kontrprzykład $w$, znajdź $s \in S$, $a \in \Sigma$ i $w'$ które jest sufiksem słowa $w$,
    takie że $S \cup \{sa\}$ i $C \cup \{w'\}$ są rozdzielne.
\end{enumerate}

Jeśli docelowy automat ma $n$ stanów, a najdłuższy zwrócony przez nauczyciela kontrprzykład $m$ liter to algorytm zadaje
$O(n)$ zapytań o równoważność i $O(mn^2)$ zapytań o należenie do języka~\cite{Angluin}.

\textbf{Uczenie się DFA z poradą.} W pracy \textit{Active Automata Learning with Advice}~\cite{srs} przedstawiono podejście,
w którym algorytm aktywnego uczenia się może otrzymać dodatkową podpowiedź, dzięki której możliwe jest zredukowanie
liczby zapytań kierowanych do nauczyciela.

Podejście to nie wymaga zmian w algorytmie ucznia. Zapytania które zadaje uczeń, nie są wprost przekazywane do nauczyciela,
tylko kierowane do systemu, który wnioskuje o nich na podstawie zadanego systemu SRS spójnego z językiem $L$.
Jeśli system nie jest w stanie odpowiedzieć na zapytanie, wtedy przekierowuje je do nauczyciela i zwraca jego odpowiedź.

Autorzy dowodzą, że jeśli $R$ jest SRS nad $\Sigma$ i $A$ jest minimalnym DFA również nad $\Sigma$,
to SRS $R$ jest spójny z $L(A)$ wtedy i tylko wtedy, gdy dla każdego stanu $q \in Q_A$ i każdej reguły
$l \to r \in R$ mamy $\delta_A(q, l) = \delta_A(q, r)$.

Ta obserwacja umożliwia przeiterowanie się po konfiguracjach hipotezy $H$ i~regułach systemu SRS $R$ w celu wykrycia
niespójności hipotezy z systemem. Jeżeli dla stanu $q \in Q_H$ oraz reguły $l \to r \in R$ zachodzi
$H(q,l) \neq H(q,r)$, to można obliczyć najkrótsze słowa $u, v \in \Sigma^*$, takie że $\delta_H(q_0, u) = q$ oraz $v$
rozróżnia stany $\delta_H(q, l)$ i $\delta_H(q, r)$. Wtedy jedno ze słów $ulv$ i $urv$ jest poprawnym kontrprzykładem,
do wybrania odpowiedniego słowa wystarczy jedno zapytanie $MQ$.

\section{Aktywne uczenie się automatów DVPA}

\textbf{Aktywne uczenie z resetami} jest rozszerzeniem modelu aktywnego uczenia zaproponowanym przez
Jana Otopa i Jakuba Michaliszyna \cite[5. Learning with resets, 6. Learning with a single reset]{main}.
Przejście z aktywnego uczenia DFA do aktywnego uczenia DVPA wymaga kilku istotnych zmian.

W tym podejściu model zakłada naukę nieznanego uczniowi automatu DVPA $T$, zwanego automatem docelowym,
zamiast języka rozpoznawanego poprzez pewien automat z tej klasy.

Ponieważ problem minimalizacji DVPA jest NP-trudny~\cite{minimization}, przy założeniu $P \neq NP$
nie istnieje wielomianowy algorytm w stylu $L^*$ konstruujący minimalny automat.
Dlatego podejście to nie gwarantuje minimalizacji automatu.

Autorzy wykazują również, że sama obserwacja stosu nie wystarczy do stworzenia algorytmu działającego
w czasie wielomianowym \cite[3.1 Difficulty of adapting $L^*$ to visibly pushdown languages]{main}.
Dlatego algorytm modyfikuje stos automatu w~trakcie wczytywania słowa.

W celu umożliwienia manipulacji zawartością stosu autorzy wprowadzają dodatkowe pojęcia związane z kontrolą stosu:
\begin{description}
    \item[Litera kontrolna] $\underline{\bot s}$, gdzie $s \in \Gamma^*$, jest to litera po przeczytaniu której automat podmienia
      zawartość stosu na $\bot s$.

    \item[Zbiór liter kontrolnych] oznaczamy symbolem $\underline{\Gamma}$.

    \item[Słowo standardowe] jest to słowo nad alfabetem $\Sigma$, nie zawiera ono liter kontrolnych.

    \item[Słowo kontrolne] jest to słowo składające się zarówno z liter z alfabetu $\Sigma$, jak również liter kontrolnych.

    \item[Słowo z pojedynczym resetem] jest to słowo zawierające dokładnie jedną literę kontrolną. Zbiór wszystkich słów z pojedynczym
      resetem nad alfabetem $\Sigma$ i $\Gamma$ oznaczamy jako $Reset_1^{\Sigma, \Gamma}$.

    \item[Słowo z pojedynczym resetem inicjalizującym] jest to słowo z pojedynczym resetem, którego pierwszą literą jest litera kontrolna.
      Zbiór tych słów oznaczamy jako $Init_1^{\Sigma,\Gamma}$. Dla uproszczenia notacji, indeks górny będzie pomijany.
\end{description}
Poprzez $\underline{L}(A)$ oznaczamy język rozpoznawany przez automat nad alfabetem zawierającym litery kontrolne.

Zbiór zapytań na jakie powinien odpowiadać nauczyciel w tym modelu zostaje rozszerzony o dodatkowe zapytanie: zapytanie o stos $SC(w)$,
w odpowiedzi na to zapytanie nauczyciel powinien zwrócić zawartość stosu automatu docelowego po przeczytaniu słowa $w$.
Zapytania zostają również dostosowane do modelu zawierającego litery kontrolne w następujący sposób:

\begin{itemize}
    \item \textbf{zapytanie o równoważność $EQ(H)$}. Nauczyciel powinien zwrócić słowo standardowe, rozróżniające język $L(H)$ od języka $L(T)$,
    gdy takie istnieje.

    \item \textbf{zapytanie o należenie $MQ(w)$}. Nauczyciel powinien odpowiedzieć czy słowo $w$ będące słowem standardowym,
    lub słowem kontrolnym należy do języka $\underline{L}(T)$.

    \item \textbf{zapytanie o stos $SC(w)$}. Nauczyciel powinien zwrócić zawartość stosu po przeczytaniu słowa $w$ będącego słowem standardowym,
    lub słowem kontrolnym.
\end{itemize}

Kluczowe jest również wprowadzenie odpowiednika relacji równoważności, która będzie udoskonalana w trakcie działania algorytmu.
Dla języków rozpoznawanych przez DVPA Twierdzenie Myhilla-Nerode'a nie może być zastosowane wprost. W szczególności minimalny automat DVPA
rozpoznający dany język nie jest jednoznaczny, a kanoniczna konstrukcja automatu~\cite{canonical} może być wykładniczo większa niż minimalny DVPA.

Autorzy wprowadzają pojęcie kontrolnej prawostronnej kongruencji. Relacja $\simeq$ na słowach nad alfabetem $(\Sigma \cup \underline{\Gamma})$,
jest kontrolną prawostronną kongruencją, jeśli:
$$\text{dla wszystkich słów }w_1, w_2 \in (\Sigma \cup \underline{\Gamma})^* \text{ i } \alpha \in \underline{\Gamma} \cdot \Sigma \text{ } (w_1 \simeq w_2 \Rightarrow w_1\alpha \simeq w_2\alpha)$$

Relacje tego typu stanowią podstawę konstrukcji automatu uczonego przez algorytm. 
W następnym rozdziale przedstawiono konkretną relację zaproponowaną przez autorów, która została zaimplementowana.

\textbf{Uczenie się DVPA z poradą.}
Podejście aktywnego uczenia z poradą można przenieść również na języki visibly pushdown.
W pracy \textit{Learning Visibly Pushdown Languages under Rewriting Rules}~\cite{SrsDvpa} zaproponowano
wykorzystanie systemów SRS dla automatów DVPA.

W przypadku języków visibly pushdown istotne jest jednak to, że słowa mają dodatkową strukturę wynikającą z dopasowania
symboli typu call i return. Dlatego nie każda reguła przepisywania jest naturalna z punktu widzenia języków visibly pushdown.
Reguły powinny zachowywać strukturę dopasowań poza przepisywanym fragmentem słowa. Dlatego rozważane są systemy
SRS zachowujące dopasowanie. W tej pracy implementacyjnie ograniczamy się do dobrze zbalansowanych systemów SRS,
które stanowią szczególny przypadek systemów zachowujących dopasowanie.

Autorzy wskazują również~\cite[Example 4]{SrsDvpa}, że niektóre naturalne własności języków visibly pushdown
nie mogą być wyrażone za pomocą skończonego systemu zwykłych reguł SRS. Z tego powodu wprowadzane
są reguły ze zmiennymi, w których za zmienne podstawiane są słowa dobrze zbalansowane. Taka reguła reprezentuje
zbiór wszystkich zwykłych reguł otrzymanych przez podstawienie za zmienną dowolnego słowa dobrze zbalansowanego.

Opis sposobu sprawdzania spójności hipotezy z systemem SRS oraz zastosowane uproszczenia implementacyjne
przedstawiono w rozdziale~\ref{chap:srs}

\section{Heurystyki wykorzystane w implementacji}

W implementacji zastosowano kilka rozwiązań o charakterze heurystycznym. Nie zmieniają one ogólnego modelu aktywnego uczenia,
ale wpływają na praktyczne działanie programu oraz interpretację wyników eksperymentalnych. Szczegóły tych
rozwiązań opisano w kolejnych rozdziałach przy omawianiu odpowiednich elementów implementacji.

\textbf{Algorytm ucznia.}
Pierwszym rozwiązaniem heurystycznym jest dodanie wymuszonych przejść.
Procedura ta została opisana w podrozdziale~\ref{chap:counterexample} W trakcie implementacji zaobserwowano,
że w~rzadkich przypadkach algorytm może zapętlać się podczas obsługi kontrprzykładów
związanych z przejściami typu return. Całkowite wyeliminowanie tego problemu wymagałoby
istotnych zmian w sposobie generowania hipotezy, które mogłyby znacząco zwiększyć koszt działania algorytmu.

Z tego powodu zastosowano heurystyczny mechanizm wymuszonych przejść. Rozwiązanie to pozwala niemal
całkowicie uniknąć zapętleń w praktycznych przypadkach testowych, wprowadzając niezauważalny narzut czasowy.

\textbf{Sprawdzanie spójności z systemem SRS.}
Drugim rozwiązaniem heurystycznym jest uproszczenie sprawdzania spójności hipotezy z systemem SRS.
Oryginalne podejście wymaga symbolicznego uwzględnienia wszystkich możliwych podstawień słów dobrze
zbalansowanych za zmienne. W implementacji zamiast tego konstruowany jest ograniczony zbiór słów
dobrze zbalansowanych. Słowa do tego zbioru wybierane są heurystycznie, tak aby premiować słowa prowadzące
w hipotezie do możliwie różnych zachowań.

Takie podejście zachowuje poprawność wykrytych niespójności. Oznacza to, że jeśli algorytm wykryje
niespójność, to hipoteza rzeczywiście nie jest zgodna z systemem SRS. Nie gwarantuje jednak wykrycia
każdej możliwej niespójności, ponieważ rozważany jest tylko ograniczony zbiór podstawień za zmienną.

\textbf{Heurystyczne sprawdzanie równoważności.}
W części eksperymentów wykorzystano również przybliżone sprawdzanie równoważności przez testowanie
losowo generowanych słów. Takie podejście może znacząco skrócić czas uczenia, ale nie gwarantuje, że
znaleziony automat jest równoważny docelowemu.